% English translation of chapter7.tex lines 712--891.
% Source: Methods of Algebra, Volume 2, commit
% 9a5803ff77dd3257484cb177f851a73770a59dd3 (CC BY 4.0).
% stable-unit-id: o014.aljabr2.chapter7.closed-structure-dgcat
% Coverage: all of Section 7.4, ``Case Study: The Closed Structure on
% dg-Categories''; ends immediately before Section 7.5 at line 892.

% segment-id: o014.li2.chapter7.closed-structure-dgcat.h001
\section{Case Study: The Closed Structure on dg-Categories}\label{sec:dg-closed}
\hypertarget{o014.aljabr2.chapter7.closed-structure-dgcat}{ }
% segment-id: o014.li2.chapter7.closed-structure-dgcat.p001
We previously discussed the closed monoidal structure on $\cate{Cat}$, where
$\otimes$ comes from the ordinary product of categories
$\mathcal{C}_1\times\mathcal{C}_2$. On the other hand, relative to a fixed
commutative ring $\Bbbk$, Definition~\ref{def:dgCat} introduced
dg-categories, namely categories enriched over
$(\cate{C}(\Bbbk\dcate{Mod}),\otimes)$. The aim of this section is to show
that all small dg-categories themselves form a closed monoidal category
$\cate{dgCat}_{\Bbbk}$. This category can be regarded as both a linearized
and a complex-enriched version of $\cate{Cat}$, and it arises naturally and
usefully in many situations. All that is needed is a sequence of lengthy but
straightforward constructions.

% segment-id: o014.li2.chapter7.closed-structure-dgcat.p002
The first step is to upgrade the ordinary product of categories to a tensor
product of dg-categories. This involves the symmetric braiding on
$(\cate{C}(\Bbbk\dcate{Mod}),\otimes)$.

% segment-id: o014.li2.chapter7.closed-structure-dgcat.d001
\begin{definition}\label{def:dgCat-tensor}
	Let $\mathcal{C}_1$ and $\mathcal{C}_2$ be dg-categories over $\Bbbk$.
	Define a new dg-category $\mathcal{C}_1\otimes\mathcal{C}_2$ by
	% segment-id: o014.li2.chapter7.closed-structure-dgcat.q001
	\begin{align*}
		\Obj(\mathcal{C}_1 \otimes \mathcal{C}_2) & := \Obj(\mathcal{C}_1) \times \Obj(\mathcal{C}_2), \\
		\Hom_{\mathcal{C}_1 \otimes \mathcal{C}_2}^\bullet((X, X'), (Y, Y')) & := \Hom_{\mathcal{C}_1}^\bullet(X, Y) \otimes \Hom_{\mathcal{C}_2}^\bullet(X', Y'),
	\end{align*}
	and take $\identity_X\otimes\identity_{X'}$ as the unit of
	$\Hom_{\mathcal{C}_1\otimes\mathcal{C}_2}^\bullet((X,X'),(X,X'))$.
	Composition of morphisms is defined by
	% segment-id: o014.li2.chapter7.closed-structure-dgcat.q002
	\[\begin{tikzcd}
		\left( \Hom_{\mathcal{C}_1}^\bullet(Y, Z) \otimes \Hom_{\mathcal{C}_2}^\bullet(Y', Z')\right) \otimes
		\left( \Hom_{\mathcal{C}_1}^\bullet(X, Y) \otimes \Hom_{\mathcal{C}_2}^\bullet(X', Y') \right) \arrow[d, "\sim" sloped] \\
		\left( \Hom_{\mathcal{C}_1}^\bullet(Y, Z) \otimes \Hom_{\mathcal{C}_1}^\bullet(X, Y) \right) \otimes
		\left( \Hom_{\mathcal{C}_2}^\bullet(Y', Z') \otimes \Hom_{\mathcal{C}_2}^\bullet(X', Y') \right) \arrow[d] \\
		\Hom_{\mathcal{C}_1}^\bullet(X, Z) \otimes \Hom_{\mathcal{C}_2}^\bullet(X', Z')
	\end{tikzcd}\]
	where the isomorphism comes from the Koszul braiding and the associativity
	constraint.
\end{definition}

% segment-id: o014.li2.chapter7.closed-structure-dgcat.p003
The symmetry of the braiding shows that
$\mathcal{C}_1\otimes\mathcal{C}_2$ is naturally equivalent to
$\mathcal{C}_2\otimes\mathcal{C}_1$. Write $\Bbbk$ for the dg-category with
one object and endomorphism complex $\Bbbk$. It is easy to prove that
$\Bbbk\otimes\mathcal{C}$ is equivalent to $\mathcal{C}\otimes\Bbbk$.

% segment-id: o014.li2.chapter7.closed-structure-dgcat.d002
\begin{definition}\label{def:dgCat-as-cat}
	\index[sym1]{dgCat@$\cate{dgCat}_{\Bbbk}$}
	Let $\Bbbk$ be a commutative ring. Write $\cate{dgCat}_{\Bbbk}$ for the
	category of all small dg-categories over $\Bbbk$. This is a symmetric
	monoidal category with unit $\Bbbk$.
\end{definition}

% segment-id: o014.li2.chapter7.closed-structure-dgcat.p004
The restriction to small categories is, of course, meant to avoid
set-theoretic difficulties. Its chief benefit is that all functors between
two given small categories $\mathcal{C}$ and $\mathcal{D}$ form a small set.

% segment-id: o014.li2.chapter7.closed-structure-dgcat.p005
The second step is to enrich the $\Hom$ set between dg-functors into a
complex.

% segment-id: o014.li2.chapter7.closed-structure-dgcat.d003
\begin{definition}
	Let $F,G:\mathcal{C}\to\mathcal{D}$ be dg-functors between dg-categories.
	For every $n\in\Z$, a morphism of degree $n$ from $F$ to $G$ is defined to
	be data
	% segment-id: o014.li2.chapter7.closed-structure-dgcat.q003
	\[ \phi = \left( \phi_X \right)_{X \in \Obj(\mathcal{C})}, \quad \phi_X \in \Hom_{\mathcal{D}}^n(FX, GX), \]
	such that, for all $m\in\Z$, $X,Y\in\Obj(\mathcal{C})$, and
	$f\in\Hom^m_{\mathcal{C}}(X,Y)$, the following equality holds in
	$\Hom^{n+m}_{\mathcal{D}}(FX,GY)$:
	% segment-id: o014.li2.chapter7.closed-structure-dgcat.q004
	\[ (Gf) \phi_X = (-1)^{nm} \phi_Y (Ff); \]
	here composition in this formula is understood as the morphisms in
	$\Bbbk\dcate{Mod}$
	% segment-id: o014.li2.chapter7.closed-structure-dgcat.q005
	\begin{gather*}
		\Hom^m_{\mathcal{D}}(GX, GY) \otimes \Hom^n_{\mathcal{D}}(FX, GX) \to \Hom^{m+n}_{\mathcal{D}}(FX, GY), \\
		\Hom^n_{\mathcal{D}}(FY, GY) \otimes \Hom^m_{\mathcal{D}}(FX, FY) \to \Hom^{m+n}_{\mathcal{D}}(FX, GY).
	\end{gather*}
	The $\Bbbk$-module of all degree-$n$ morphisms $\phi:F\to G$ is denoted by
	$\iHom^n(F,G)$.
\end{definition}

% segment-id: o014.li2.chapter7.closed-structure-dgcat.p006
The condition above may be viewed as saying that the diagram of graded
morphisms
% segment-id: o014.li2.chapter7.closed-structure-dgcat.q006
\[\begin{tikzcd}
	FX \arrow[r, "{\phi_X}"] \arrow[d, "Ff"'] & GX \arrow[d, "Gf"] \\
	FY \arrow[r, "{\phi_Y}"'] & GY
\end{tikzcd} \quad
\begin{array}{cc}
	\phi: \text{degree }n \\
	f: \text{degree }m
\end{array}\]
commutes up to the sign $(-1)^{nm}$ required by the Koszul sign rule.
Strictly speaking, these ``graded morphisms'' are not morphisms, but must be
understood as elements of a $\Hom$ complex.

% segment-id: o014.li2.chapter7.closed-structure-dgcat.dp001
\begin{definition-proposition}
	For dg-functors $F,G:\mathcal{C}\to\mathcal{D}$ and every $n\in\Z$, one
	can define a homomorphism
	% segment-id: o014.li2.chapter7.closed-structure-dgcat.q007
	\[ d^n = d^n_{\iHom^\bullet(F, G)}: \iHom^n(F, G) \to \iHom^{n+1}(F, G). \]

	For every $\phi=(\phi_X)_{X\in\Obj(\mathcal{C})}$, set
	% segment-id: o014.li2.chapter7.closed-structure-dgcat.q008
	\[ (d^n \phi)_X := d^n_{\Hom^\bullet(FX, GX)}(\phi_X). \]
	This makes $\left(\iHom^n(F,G),d^n\right)_{n\in\Z}$ into a complex
	$\iHom^\bullet(F,G)$.
\end{definition-proposition}
% segment-id: o014.li2.chapter7.closed-structure-dgcat.pf001
\begin{proof}
	First check that $d^n\phi$ does indeed belong to
	$\iHom^{n+1}(F,G)$. Let $f\in\Hom^m_{\mathcal{C}}(X,Y)$. Apply
	$d^{m+n}$ to both sides of
	$(Gf)\phi_X=(-1)^{nm}\phi_Y(Ff)$ (suppressing the subscript
	$\Hom^\bullet(FX,GY)$). This gives
	% segment-id: o014.li2.chapter7.closed-structure-dgcat.q009
	\begin{align*}
		d^{m+n}((Gf) \phi_X) & = (d^m Gf) \phi_X + (-1)^m (Gf) (d^n \phi_X) \\
		& = G (d^m f) \phi_X + (-1)^m (Gf) (d^n \phi)_X \\
		& = (-1)^{n(m+1)} \phi_Y F(d^m f) + (-1)^m (Gf) (d^n \phi)_X, \\
		(-1)^{nm} d^{m+n}(\phi_Y (Ff)) & = (-1)^{nm} (d^n \phi_Y) Ff + (-1)^{nm+n} \phi_Y d^n(Ff) \\
		& = (-1)^{nm} (d^n \phi)_Y Ff + (-1)^{nm+n} \phi_Y F(d^m f).
	\end{align*}
	Equality of these two expressions immediately yields
	$(Gf)(d^n\phi)_X=(-1)^{(n+1)m}(d^n\phi)_YFf$.

	Second, since $\Hom^\bullet(FX,GX)$ is a complex,
	% segment-id: o014.li2.chapter7.closed-structure-dgcat.q010
	\begin{align*}
		(d^{n+1} d^n \phi)_X & = d^{n+1}_{\Hom^\bullet(FX, GX)}((d^n \phi)_X) \\
		& = d^{n+1}_{\Hom^\bullet(FX, GX)} d^n_{\Hom^\bullet(FX, GX)} \phi_X = 0.
	\end{align*}
	Thus we obtain the complex $\iHom^\bullet(F,G)$.
\end{proof}

% segment-id: o014.li2.chapter7.closed-structure-dgcat.p007
For three dg-functors $F,G,H:\mathcal{C}\to\mathcal{D}$, there is a
composition operation
% segment-id: o014.li2.chapter7.closed-structure-dgcat.q011
\begin{equation}\label{eqn:dgfct-iHom}\begin{aligned}
	\iHom^a(G, H) \otimes \iHom^b(F, G) & \to \iHom^{a+b}(F, H) \\
	\psi \otimes \phi & \mapsto \psi\phi := \left( \psi_X \phi_X \right)_{X \in \Obj(\mathcal{C})},
\end{aligned}\end{equation}
where $a,b\in\Z$. This operation is associative and satisfies
$d^{a+b}(\psi\phi)=(d^a\psi)\phi+(-1)^a\psi(d^b\phi)$. By definition,
everything reduces to a verification in $\Hom^\bullet_{\mathcal{D}}$. This
composition should be understood as vertical composition of graded
morphisms, depicted as
% segment-id: o014.li2.chapter7.closed-structure-dgcat.q012
\[\begin{tikzcd}
	\mathcal{C}
	\arrow[bend left=70, rr, "F", ""' name=U]
	\arrow[rr, "G" name=MM, ""' name=M]
	\arrow[bend right=70, rr, "" name=D, "H"'] & &
	\arrow[Rightarrow, to path=(U) -- (MM) \tikztonodes, "\phi"] \arrow[Rightarrow, to path=(M) -- (D) \tikztonodes, "\psi"] \mathcal{D}
\end{tikzcd}  \quad \text{composed to give} \quad \begin{tikzcd}
	\mathcal{C}
	\arrow[bend left=50, rr, "F", ""' name=U]
	\arrow[bend right=50, rr, "" name=D, "H"']
	& & \arrow[Rightarrow, to path=(U) -- (D) \tikztonodes, "\psi \phi"] \mathcal{D} .
\end{tikzcd}\]

% segment-id: o014.li2.chapter7.closed-structure-dgcat.p008
Consider the special case $n=0$. The condition on
$\phi=(\phi_X)_X\in\iHom^0(F,G)$ says precisely that
$(Gf)\phi_X=\phi_Y(Ff)$ for every $f\in\Hom^m(X,Y)$, while the condition
$d^0\phi=0$ says precisely that
$\phi_X\in\Ker\left(d^0_{\Hom^\bullet(FX,GX)}\right)$ for every
$X\in\Obj(\mathcal{C})$. Hence the elements of the $\Bbbk$-module
% segment-id: o014.li2.chapter7.closed-structure-dgcat.q013
\[ \Ker\left[ \iHom^0(F, G) \xrightarrow{d^0} \iHom^1(F, G) \right] \]
are exactly the morphisms $\phi:F\to G$ between the dg-functors $F$ and $G$
in the ordinary sense.

% segment-id: o014.li2.chapter7.closed-structure-dgcat.p009
The next step is to enrich the functor category between dg-categories into a
dg-category.

% segment-id: o014.li2.chapter7.closed-structure-dgcat.d004
\begin{definition}
	For arbitrary dg-categories $\mathcal{C}$ and $\mathcal{D}$, define the
	dg-category $\iHom(\mathcal{C},\mathcal{D})$ as follows.
	\begin{itemize}
		\item Its objects are all dg-functors $F:\mathcal{C}\to\mathcal{D}$.
		\item For any dg-functors $F,G:\mathcal{C}\to\mathcal{D}$, their
		$\Hom$ complex is the complex $\iHom^\bullet(F,G)$ defined above.
		\item The composition operation
		$\iHom^\bullet(G,H)\otimes\iHom^\bullet(F,G)\to
		\iHom^\bullet(F,H)$ is defined as in \eqref{eqn:dgfct-iHom}; it is
		called vertical composition.
	\end{itemize}
\end{definition}

% segment-id: o014.li2.chapter7.closed-structure-dgcat.p010
If $\mathcal{C}$ and $\mathcal{D}$ are both small dg-categories, then
$\iHom(\mathcal{C},\mathcal{D})$ is also small: its set of objects is a small
set.

% segment-id: o014.li2.chapter7.closed-structure-dgcat.p011
In Definition~\ref{def:dgCat-tensor} we defined the tensor product of two
dg-categories, making the category $\cate{dgCat}_{\Bbbk}$ of
Definition~\ref{def:dgCat-as-cat} a symmetric monoidal category under
$\otimes$. The relation between $\otimes$ and $\iHom$ can now be stated
clearly as follows.

% segment-id: o014.li2.chapter7.closed-structure-dgcat.pr001
\begin{proposition}\label{prop:dgCat-closed}
	The symmetric monoidal category $\cate{dgCat}_{\Bbbk}$ is closed. Its
	internal Hom is given by $\iHom(\mathcal{C},\mathcal{D})$ as defined above,
	where $\mathcal{C}$ and $\mathcal{D}$ range over all small dg-categories.
\end{proposition}
% segment-id: o014.li2.chapter7.closed-structure-dgcat.pf002
\begin{proof}
	The key is to give a canonical isomorphism
	% segment-id: o014.li2.chapter7.closed-structure-dgcat.q014
	\[ \Hom_{\cate{dgCat}_{\Bbbk}}(\mathcal{C} \otimes \mathcal{D}, \mathcal{E}) \simeq \Hom_{\cate{dgCat}_{\Bbbk}}(\mathcal{C}, \iHom(\mathcal{D}, \mathcal{E}) ). \]

	If the dg-structure is ignored, the matter is simple: specifying
	$F:\mathcal{C}\times\mathcal{D}\to\mathcal{E}$ is equivalent to
	specifying a functor $\mathcal{C}\to\mathcal{E}^{\mathcal{D}}$ that maps
	an object $X$ to the functor
	$F(X,\mathord\cdot):\mathcal{D}\to\mathcal{E}$. The essential issue is
	the condition at the level of $\Hom$ that makes a functor a dg-functor.

	By definition, making $F$ into a dg-functor is equivalent to upgrading the
	maps on $\Hom$ sets to morphisms of complexes
	% segment-id: o014.li2.chapter7.closed-structure-dgcat.q015
	\[ F_{X, Y, Z, W}: \Hom^\bullet_{\mathcal{C}}(X, Y) \otimes \Hom^\bullet_{\mathcal{D}}(Z, W) \to \Hom^\bullet_{\mathcal{E}}(F(X, Z), F(Y, W)), \]
	functorial in all four variables. Treating the two variables of $F$
	separately, these data can be split into two parts:
	\begin{itemize}
		\item for every $X\in\Obj(\mathcal{C})$, the functor
		$F(X,\mathord\cdot)$ is upgraded to a dg-functor;
		\item there is a family of morphisms
		$F_{X,Y,Z}:\Hom^\bullet_{\mathcal{C}}(X,Y)\to
		\Hom^\bullet_{\mathcal{E}}(F(X,Z),F(Y,Z))$, functorial in all three
		variables.
	\end{itemize}
	For illustration, $F_{X,Y,Z,W}(f\otimes g)$ is obtained from
	$F_{X,Y,Z}(f)$ by using $g\in\Hom^b_{\mathcal{D}}(Z,W)$ through
	$F(Y,\mathord\cdot)$. The data in the second item are equivalent to a
	family of morphisms
	% segment-id: o014.li2.chapter7.closed-structure-dgcat.q016
	\[ \Hom^\bullet_{\mathcal{C}}(X, Y) \to \iHom^\bullet(F(X, \cdot), F(Y, \cdot)), \]
	functorial in $X$ and $Y$. In conclusion, specifying the dg-functor $F$
	is equivalent to specifying a dg-functor from $\mathcal{C}$ to
	$\iHom(\mathcal{D},\mathcal{E})$.
\end{proof}

% segment-id: o014.li2.chapter7.closed-structure-dgcat.p012
Together with the theory in \S\ref{sec:closed-monoidal}, the preceding
discussion shows that $\cate{dgCat}_{\Bbbk}$ is enriched over itself: for any
small dg-categories $\mathcal{C}$ and $\mathcal{D}$, their $\iHom$ object is
again a dg-category.

% segment-id: o014.li2.chapter7.closed-structure-dgcat.p013
Since $\otimes$ and $\iHom$ determine each other through the adjunction
\eqref{eqn:closed-monoidal-adj0}, once either the definition of $\otimes$ or
that of $\iHom$ is accepted, the other definition is at least equally
natural.

% segment-id: o014.li2.chapter7.closed-structure-dgcat.r001
\begin{remark}
	With Proposition~\ref{prop:dgCat-closed} in hand, the general properties of
	closed monoidal categories show that for any three small dg-categories
	$\mathcal{C}$, $\mathcal{D}$, and $\mathcal{E}$, composition of functors
	can be upgraded to a dg-functor
	% segment-id: o014.li2.chapter7.closed-structure-dgcat.q017
	\[ \iHom(\mathcal{D}, \mathcal{E}) \otimes \iHom(\mathcal{C}, \mathcal{D}) \to \iHom(\mathcal{C}, \mathcal{E}). \]
	On objects, this functor maps $(G,F)$ to $GF$; on $\Hom$ complexes, it is
	written
	% segment-id: o014.li2.chapter7.closed-structure-dgcat.q018
	\begin{equation}\begin{gathered}\label{eqn:dgCat-fct-compose}
		\begin{tikzcd}[row sep=small, column sep=small]
			\iHom^\bullet(G_1, G_2) \otimes \iHom^\bullet(F_1, F_2) \arrow[r] \arrow[equal, d] & \iHom^\bullet(G_1 F_1, G_2 F_2) \\
			\Hom^\bullet_{ \iHom(\mathcal{D}, \mathcal{E}) \otimes \iHom(\mathcal{C}, \mathcal{D}) }((G_1, F_1), (G_2, F_2)) &
		\end{tikzcd} \\
		F_i: \mathcal{C} \to \mathcal{D}, \quad G_i: \mathcal{D} \to \mathcal{E}, \quad i = 1, 2 .
	\end{gathered}\end{equation}

	The operation above may be pictured as horizontal composition of
	morphisms between dg-functors, with diagram
	% segment-id: o014.li2.chapter7.closed-structure-dgcat.q019
	\[ \begin{tikzcd}
		\mathcal{C} \arrow[bend left=50, r, "F_1", ""' name=LU] \arrow[bend right=50, r, "" name=LD, "F_2"'] &
		\mathcal{D} \arrow[bend left=50, r, "G_1", ""' name=RU] \arrow[bend right=50, r, "" name=RD, "G_2"'] &
		\arrow[Rightarrow, to path=(LU) -- (LD) \tikztonodes, "\phi"] \arrow[Rightarrow, to path=(RU) -- (RD) \tikztonodes, "\psi"] \mathcal{E}
	\end{tikzcd} \quad \text{composed to give} \quad \begin{tikzcd}
		\mathcal{C}
		\arrow[bend left=50, rr, "G_1 F_1", ""' name=U]
		\arrow[bend right=50, rr, "" name=D, "G_2 F_2"']
		& & \arrow[Rightarrow, to path=(U) -- (D) \tikztonodes, "\psi \phi"] \mathcal{E} .
	\end{tikzcd} \]
	Because this gives a functor, \eqref{eqn:dgCat-fct-compose} must
	\begin{inparaenum}[(a)]
		\item be associative,
		\item commute with pullback (or pushforward) along graded morphisms in
		the dg-category $\iHom(\mathcal{D},\mathcal{E})$ (or
		$\iHom(\mathcal{C},\mathcal{D})$).
	\end{inparaenum}
	But the operation in (b) is precisely the vertical composition introduced
	in \eqref{eqn:dgfct-iHom}. Thus, we see that the order of vertical and
	horizontal composition can be interchanged.

	For functors between ordinary categories, morphisms between the functors
	likewise have commuting vertical and horizontal compositions. This can be
	understood as a property of the Cartesian closed category $\cate{Cat}$ (or
	by regarding it as a $2$-category; see \cite[Example 3.5.3]{Li1}). The
	property of $\cate{dgCat}_{\Bbbk}$ above is its dg version, or its linear
	algebra version.
\end{remark}
